% Generated draft: requires mathematical and source review.
\documentclass[11pt]{article}
\usepackage[margin=25mm]{geometry}
\usepackage{fontspec}
\setmainfont{Latin Modern Roman}
\usepackage{amsmath,amssymb,mathtools,unicode-math}
\setmathfont{Latin Modern Math}
\usepackage{enumitem,needspace,xurl,hyperref}
\setlength{\parindent}{0pt}
\setlength{\parskip}{5pt}
\setlength{\emergencystretch}{3em}
\newcommand{\N}{\mathbb{N}}
\newcommand{\Z}{\mathbb{Z}}
\newcommand{\Q}{\mathbb{Q}}
\newcommand{\R}{\mathbb{R}}
\newcommand{\C}{\mathbb{C}}
\newcommand{\F}{\mathbb{F}}
\newcommand{\A}{\mathbb{A}}
\newcommand{\PP}{\mathbb{P}}
\newcommand{\E}{\mathbb{E}}
\newcommand{\eps}{\varepsilon}
\newcommand{\dd}{\,\mathrm{d}}
\begin{document}
\section*{UnsolvedMath ID: 701716}
\textbf{Source problem code:} catalogue entry
\subsection{Definitions and mathematical statement}
For a finite group \(G\), its spectrum is
\[
\omega(G)=\{\operatorname{ord}(g):g\in G\},
\]
where \(\operatorname{ord}(g)\) is the order of the element \(g\). The group \(G\) is recognizable by spectrum if, for every finite group \(H\),
\[
\omega(H)=\omega(G)\quad\Longrightarrow\quad H\cong G.
\]
The socle \(\operatorname{Soc}(G)\) is the subgroup generated by all minimal nontrivial normal subgroups of \(G\). A nonabelian simple group is a nonabelian group whose only normal subgroups are the identity subgroup and the whole group. For a group \(L\) and a positive integer \(k\), write \(L^k=\prod_{i=1}^{k}L\) for its \(k\)-fold direct power.

The source asks the following two questions. The indices \(n\) and \(k\) are interpreted as positive integers, as appropriate for the direct powers in the questions; their ranges are not explicitly stated in the source.
\begin{enumerate}
\item[(a)] Is it true that for every positive integer \(n\) there exists a finite nonabelian simple group \(S\), allowed to depend on \(n\), such that \(S^n\) is recognizable by spectrum? Equivalently, is
\[
\forall n\in\Z_{\geq 1}\ \exists S\text{ finite nonabelian simple}\ \forall H\text{ finite},\qquad
\omega(H)=\omega(S^n)\ \Longrightarrow\ H\cong S^n
\]
true?
\item[(b)] Is there a finite nonabelian simple group \(L\) such that for every positive integer \(n\) there exist an integer \(k\geq n\) and a finite group \(G\) satisfying
\[
\operatorname{Soc}(G)\cong L^k
\quad\text{and}\quad
\forall H\text{ finite},\qquad
\omega(H)=\omega(G)\ \Longrightarrow\ H\cong G?
\]
Here \(L\) is fixed independently of \(n\), while \(k\) and \(G\) may depend on \(n\). The finiteness of \(L\) follows from its occurrence as a factor of the socle of a finite group with \(k\geq1\).
\end{enumerate}
The supplied record labels the overall problem partially solved and reports an affirmative solution to part (a). These are supplied source assessments, not independently verified mathematical conclusions. No separate resolution of part (b) is supplied.
\subsection{Short English statement}
Part (a) asks whether every positive direct-power exponent occurs for a finite nonabelian simple group whose direct power is determined among finite groups by its element orders. Part (b) asks whether one fixed simple group can occur with unbounded multiplicity in the socles of finite groups determined by their element orders. The supplied source reports part (a) solved affirmatively and labels the overall problem partially solved.
\subsection{Sources}
\begin{itemize}
\item Local UnsolvedMath registry, permanent ID \(701716\), source code \texttt{catalogue entry}, title \emph{Kourovka Notebook Problem 20.58}; supplied complete source record, upstream revision \texttt{d371ee917ac9647fbb06cd6fc529deec6f3dda82}.
\item V. D. Mazurov and A. V. Vasil'ev, proposers, 2022. The supplied record attributes the problem to \emph{The Kourovka Notebook (Unsolved Problems in Group Theory)}, 20th issue (2022), problem 20.58, and also identifies the 21st issue edition (2026), \texttt{arXiv:1401.0300v48}. The collection description names E. I. Khukhro and V. D. Mazurov as editors.
\item For the reported affirmative solution of part (a), the supplied discussion cites N. Yang, I. Gorshkov, A. Staroletov, and A. V. Vasil'ev, \emph{Annali Matem. Pura Appl.} \textbf{202} (2023), 2699--2714. No article title, DOI, or URL is supplied.
\item The supplied status review is dated \texttt{2026-10-16}, identifies its evidence as a maintained source, and reports that the 21st issue (October 2026), \texttt{arXiv:1401.0300v48}, marks part (a) solved affirmatively with the preceding citation. This is supplied dated curation, not an independently verified literature assessment. The registry status is \texttt{partially\_solved}.
\end{itemize}
\end{document}
